\documentclass{book}
\usepackage{lscape}  
 

\title{A Simple LaTeX Book Example}
\author{Harshvardhan Pujari}
\date{\today}

\begin{document}
	
	\maketitle  
	
	\tableofcontents  
	\newpage
	
	\chapter{Introduction}

	
	LaTeX is a typesetting system used for producing high-quality documents\footnote{LaTeX was originally developed by Leslie Lamport in the early 1980s.}. It separates content from style, helping writers focus on structure and logic rather than design.
	
	\section{Why Use LaTeX?}
	
	LaTeX automatically formats things like citations, references, and figures. For example, this very sentence is on page~\thepage.
	
	\newpage  
	
	\chapter{Landscape Example}
	
	Sometimes you might want a wide table or image that fits better on a landscape page. You can switch to landscape like this:
	
	\begin{landscape}
		\section{Wide Table}
		
		\begin{tabular}{|c|c|c|c|}
			\hline
			Chapter & Title & Pages & Notes \\
			\hline
			1 & Introduction & 1-5 & Basics of LaTeX \\
			2 & Getting Started & 6-12 & Sample syntax \\
			3 & Landscape Example & 13-15 & Wide tables/images \\
			4 & Conclusion & 16 & Wrap-up thoughts \\
			\hline
		\end{tabular}
	\end{landscape}
	
	\newpage 
	
	\chapter{Conclusion}
	
	Using LaTeX for writing books makes it easy to handle:
	\begin{itemize}
		\item Automatic Table of Contents
		\item Page numbering
		\item Cross-references
		\item Footnotes\footnote{This is another footnote, useful for explanations.}
	\end{itemize}
	
	And as you’ve seen, you can also create landscape pages for wide content.
	
\end{document}
